% -*- coding: utf-8 -*-
% !TEX program = xelatex
\documentclass[plain,noanswer,medmath]{../../template/jnuexam}
\usepackage{../../template/kysx}

\begin{document}


\examtitle{1992}{4}


\makepart{填空题}{本题共5小题，每小题3分，满分15分}


\begin{problem}
设商品的需求函数为$Q = 100 - 5P$，其中$Q,P$分别表示为需求量和价格，如果商品需求弹性的绝对值大于$1$，则商品价格的取值范围是 \fillin{}．
\end{problem}


\begin{problem}
级数$\sum_{n = 1}^{\infty}\frac{(x - 2)^{2n}}{n4^n}$的收敛域为 \fillin{}．
\end{problem}


\begin{problem}
交换积分次序$\int_0^1\dy \int_{\sqrt{y}}^{\sqrt{2 - y^2}} f(x,y)\dx =$ \fillin{}．
\end{problem}


\begin{problem}
设$A$为$m$阶方阵，$B$为$n$阶方阵，且$|A| = a$, $|B| = b$, $C = \begin{pmatrix}
  O & A \\
  B & O
\end{pmatrix}$，则$|C| =$ \fillin{}．
\end{problem}


\begin{problem}
将 C, C, E, E, I, N, S 这七个字母随机地排成一行，那么恰好排成英文单词~\hbox{SCIENCE}~的概率为 \fillin{}．
\end{problem}


\makepart{选择题}{本题共5小题，每小题3分，满分15分}


\begin{problem}
设$F(x) = \frac{x^2}{x - a}\int_a^x f(t) \dt$，其中$f(x)$为连续函数，则$\lim_{x \to a} F(x)$等于\pickout{}
\begin{abcd}
\item $a^2$．        
\item $a^2 f(a)$．        
\item $0$．        
\item 不存在．
\end{abcd}
\end{problem}


\begin{problem}
当$x \to 0$时，下面四个无穷小量中，哪一个是比其他三个更高阶的无穷小量？\pickout{}
\begin{abcd}
\item $x^2$．     
\item $1 - \cos x$．     
\item $\sqrt{1 - x^2} - 1$．     
\item $x - \tan x$．
\end{abcd}
\end{problem}


\begin{problem}
设$A$为$m \times n$矩阵，齐次线性方程组$Ax = 0$仅有零解的充分条件是\pickout{}
\begin{abcd}
\item $A$的列向量线性无关．               
\item $A$的列向量线性相关．
\item $A$的行向量线性无关．               
\item $A$的行向量线性相关．
\end{abcd}
\end{problem}


\begin{problem}
设当事件$A$与$B$同时发生时，事件$C$必发生，则\pickout{}
\begin{abcd}
\item $P(C) \le P(A) + P(B) - 1$．
\item $P(C) \ge P(A) + P(B)- 1$．
\item $P(C)= P(AB)$．
\item $P(C) = P(A \cup B)$．
\end{abcd}
\end{problem}


\begin{problem}
设$n$个随机变量$X_1$, $X_2$, $\cdots$, $X_n$独立同分布，
$$D(X_1) = \sigma^2,\quad \overline{X}= \frac{1}{n}\sum_{i = 1}^n X_i,\quad
  S^2 = \frac{1}{n - 1}\sum_{i = 1}^n (X_i - \overline{X})^2,$$
则\pickout{}
\begin{abcd}
\item $S$是$\sigma$的无偏估计量．					
\item $S$是$\sigma$的最大似然估计量．
\item $S$是$\sigma$的相合估计量（即一致估计量）．	
\item $S$与$\overline{X}$相互独立．
\end{abcd}
\end{problem}


\makepart{}{本题满分5分}

\begin{problem}
设函数$f(x) = \begin{cases}
  \frac{\ln\cos(x-1)}{1-\sin\frac{\pi x}{2}}, & x \ne 1, \\
                                           1, & x = 1.
\end{cases}$问函数$f(x)$在$x = 1$处是否连续？若不连续，修改函数在$x = 1$处的定义使之连续．
\end{problem}


\makepart{}{本题满分5分}

\begin{problem}
计算$I = \int \frac{\arccot\e^x}{\e^x} \dx$．
\end{problem}


\makepart{}{本题满分5分}

\begin{problem}
设$z = \sin(xy) + \varphi(x,\frac{x}{y})$，求$\frac{\pd^2 z}{\pdx\pd y}$，其中$\varphi (u,v)$有二阶偏导数．
\end{problem}


\makepart{}{本题满分5分}

\begin{problem}
求连续函数$f(x)$，使它满足$f(x) + 2\int_0^x f(t)\dt = x^2$．
\end{problem}


\makepart{}{本题满分6分}

\begin{problem}
求证：当$x \ge 1$时，$\arctan x - \frac{1}{2}\arccos \frac{2x}{1 + x^2} = \frac{\pi}{4}$．
\end{problem}


\makepart{}{本题满分9分}

\begin{problem}
设曲线方程$y = \e^{- x}$（$x \ge 0$）．
\begin{enumerate}
\item 把曲线$y = \e^{- x}$,$x$轴，$y$轴和直线$x = \xi(\xi > 0)$所围成平面图形绕$x$轴旋转一周，
得一旋转体，求此旋转体体积$V(\xi)$；求满足$V(a) = \frac{1}{2}\lim_{\xi \to +\infty} V(\xi)$的$a$．
\item 在此曲线上找一点，使过该点的切线与两个坐标轴所夹平面图形的面积最大，并求出该面积．
\end{enumerate}
\end{problem}


\makepart{}{本题满分7分}

\begin{problem}
设矩阵$A$与$B$相似，其中
$$A = \begin{pmatrix}
  -2 & 0 & 0 \\
   2 & x & 2 \\
   3 & 1 & 1
\end{pmatrix},\qquad B = \begin{pmatrix}
  -1 & 0 & 0 \\
   0 & 2 & 0 \\
   0 & 0 & y
\end{pmatrix}.$$
\begin{enumerate*}
\item 求$x$和$y$的值．
\item 求可逆矩阵$P$，使得$P^{-1} AP = B$．
\end{enumerate*}
\end{problem}


\makepart{}{本题满分6分}

\begin{problem}
已知三阶矩阵$B \ne 0$，且$B$的每一个列向量都是以下方程组的解:
$$\begin{cases}
  x_1 + 2x_2 - 2x_3 = 0, \\
  2x_1 - x_2 + \lambda x_3 = 0, \\
  3x_1 + x_2 - x_3 = 0.
\end{cases}$$
\begin{enumerate*}
\item 求$\lambda$的值；
\item 证明$|B| = 0$．
\end{enumerate*}
\end{problem}


\makepart{}{本题满分6分}

\begin{problem}
设$A$,$B$分别为$m$,$n$阶正定矩阵，试判定分块矩阵$C = \begin{pmatrix}
  A & 0 \\
  0 & B
\end{pmatrix}$是否是正定矩阵．
\end{problem}


\makepart{}{本题满分7分}

\begin{problem}
假设测量的随机误差$X\sim N(0,10^2)$，试求$100$次独立重复测量中，至少有三次测量误差的绝对值大于$19.6$的概率$\alpha$，并利用泊松分布求出$\alpha$的近似值（要求小数点后取两位有效数字）．
\par\centerline{[附表]}\noindent
$$\begin{array}{|c|*{8}{c}|}
\hline
        \lambda &     1 &     2 &     3 &     4 &     5 &     6 &     7 & \cdots \\
\hline
  \e^{-\lambda} & 0.368 & 0.135 & 0.050 & 0.018 & 0.007 & 0.002 & 0.001 & \cdots \\
\hline
\end{array}$$
\end{problem}


\makepart{}{本题满分5分}

\begin{problem}
一台设备由三大部分构成，在设备运转中各部件需要调整的概率相应为$0.10$, $0.20$ 和$0.30$．
假设各部件的状态相互独立，以$X$表示同时需要调整的部件数，试求$X$的概率分布，数学期望$EX$和方差$DX$．
\end{problem}


\makepart{}{本题满分4分}

\begin{problem}
设二维随机变量$(X,Y)$的概率密度为
$$f(x,y) = \begin{cases}
  \e^{- y}, & 0 < x < y, \\
         0, & \text{其他},
\end{cases}$$
\begin{enumerate*}
\item 求随机变量$X$的密度$f_X(x)$；
\item 求概率$P\{X + Y \le 1\}$．
\end{enumerate*}
\end{problem}


\end{document}
